% ============================================================
% AppendixA.tex — Supplementary Calculations
% ============================================================
\chapter{Supplementary Calculations}
\label{app:calculations}

\section{Cost Normalization}

For supplementary reference, the historical purchase cost of the Armfield CEK-A may be adjusted using the Chemical Engineering Plant Cost Index (CEPCI)~\cite{towler2013}:

\begin{equation}
C_{\Theta_2} = C_{\Theta_1} \times \frac{CI_{\Theta_2}}{CI_{\Theta_1}} = 525{,}700 \times \frac{797.9}{381.1} = \text{PHP } 1{,}100{,}645.58
\label{eq:inflation}
\end{equation}

where $\Theta_1 = 1995$, $CI_{\Theta_1} = 381.1$, $\Theta_2 = 2023$, and $CI_{\Theta_2} = 797.9$.

\section{Physical Design and Motor-Sizing Basis for MIXR-1}

The physical vessel geometry, impeller dimensions, and agitation-related sizing used in MIXR-1 follow the study by Dumalagan~\cite{dumalagan2024}. These calculations are placed in this appendix so that the main methodology in Section~V can remain focused on the Computer Engineering (CpE) implementation.

\subsection{Geometric Scaling}

The scaling ratio between the Armfield CEK-A and MIXR-1 is:

\begin{equation}
R = \frac{D_2}{D_1} = \frac{200}{288} = 0.6944
\label{eq:scaling_ratio}
\end{equation}

\noindent where $D_1 = 288$\,mm is the CEK-A tank inner diameter and $D_2 = 200$\,mm is the MIXR-1 tank inner diameter. All linear dimensions are scaled using:

\begin{equation}
X_2 = R \times X_1
\label{eq:scaling_general}
\end{equation}

This preserves the key dimensionless geometric ratios used for the agitation experiments and for comparison with the Brown et~al.~\cite{brown1950} correlations.

\subsection{Tank Volume and Volumetric Power Estimate}

Using the scaled tank dimensions, the working vessel volume is:

\begin{equation}
V = \frac{\pi}{4} D_t^2 H_t = \frac{\pi}{4} (0.200)^2 (0.300) = 9.42 \times 10^{-3} \text{ m}^3
\label{eq:tank_volume}
\end{equation}

Using the rule-of-thumb agitation power intensity $Q = 1.5$\,kW/m$^3$~\cite{harnby2000}, the baseline power estimate is:

\begin{equation}
P_0 = Q \times V = 1500 \times 9.42 \times 10^{-3} = 14.14 \text{ W}
\label{eq:power_volumetric}
\end{equation}

Applying a 1.5$\times$ design margin gives:

\begin{equation}
P = 1.5 \times P_0 = 1.5 \times 14.14 = 21.21 \text{ W}
\label{eq:power_with_margin}
\end{equation}

\subsection{Speed and Torque Requirement}

For the most demanding impeller configuration, the required speed from the power-number relationship is:

\begin{equation}
n = \left( \frac{P_0}{\rho \, N_{Po} \, D_i^5} \right)^{1/3} = \left( \frac{14.14}{998 \times 5.46 \times (0.12)^5} \right)^{1/3}
\label{eq:speed_calc}
\end{equation}

which evaluates to:

\begin{equation}
n = \left( \frac{14.14}{0.1356} \right)^{1/3} = (104.3)^{1/3} = 4.70 \text{ rev/s} = 282 \text{ RPM}
\label{eq:speed_result}
\end{equation}

The corresponding torque requirement is:

\begin{equation}
\tau = \frac{P}{\omega} = \frac{21.21}{2\pi \times 4.70} = \frac{21.21}{29.53} = 0.72 \text{ N}\cdot\text{m}
\label{eq:torque}
\end{equation}

These appendix calculations define the mechanical operating envelope that the CpE instrumentation and control subsystems must support.